\chapter{October 1917}

The Russian revolution of 1917 was a turning-point in history, and may well be assessed by future historians as the greatest event of the twentieth century. Like the French revolution, it will continue for a long time to polarize opinion, being hailed by some as a landmark in the emancipation of mankind from past oppression, and denounced by others as a crime and a disaster. It presented the first open challenge to the capitalist system, which had reached its peak in Europe at the end of the nineteenth century. Its occurrence at the height of the first world war, and partly as a result of that war, was more than a coincidence. The war had struck a deadly blow at the international capitalist order as it had existed before 1914, and revealed its inherent instability. The revolution may be thought of both as a consequence and as a cause of the decline of capitalism. 

While, however, the revolution of 1917 had a world-wide significance, it was also rooted in specifically Russian conditions. The imposing facade of the Tsarist autocracy concealed a stagnant rural economy, which had made few substantial advances since the emancipation of the serfs, and a hungry and restive peasantry. Terrorist groups had been at work since the eighteen-sixties with recurrent outbreaks of violence and repression. This period saw the rise of the narodnik movement, later succeeded by the Social-Revolutionary Party, whose appeal was to the peasants. After 1890 industrialization began to make important inroads into the primitive Russian economy; and the rise of an increasingly influential and wealthy industrial and financial class, heavily dependent on foreign capital, encouraged the infiltration of some liberal western ideas, which found their fullest expression in the Kadet (Constitutional-Democratic) Party. But this process was accompanied by the growth of a proletariat of factory workers and by the early symptoms of proletarian turbulence; the first strikes occurred in the eighteen-nineties. 
% 注：修复了原文 "(a stagnant" 中多余的括号。
These developments were reflected in the foundation in 1898 of the Marxist Russian Social-Democratic Workers' Party, the party of Lenin, Martov and Plekhanov. Seething unrest was brought to the surface by the frustrations and humiliations of the Russo-Japanese war. 
% 注：此处原文因为多栏扫描导致严重串行错乱 ("...reflected tnov. Seething unrest was brought to the surface by the in the foundation in 1898... Plekh- rustrations...")，已根据上下文完全恢复正确的语义逻辑。

The first Russian revolution of 1905 had a mixed character. It was a revolt of bourgeois liberals and constitutionalists against an arbitrary and antiquated autocracy. It was a revolt of workers, sparked off by the atrocity of ``Bloody Sunday'', and leading to the election of the first Petersburg Soviet of Workers' Deputies. It was a widespread revolt of peasants, spontaneous and uncoordinated, often extremely bitter and violent. The three strands were never woven together, and the revolution was easily put down at the cost of some largely unreal constitutional concessions. 
% 注：清理了原 OCR 中夹杂的乱码 " \textbullet{}*\textasciicircum{}\textbar{} " 以及 "\textasciicircum{}Fp I"，并修复了严重的句子错位重组问题。

The same factors inspired the revolution of February 1917, but this time were reinforced and dominated by war-weariness and universal discontent with the conduct of the war. Nothing short of the Tsar's abdication could stem the tide of revolt. Autocracy was replaced by the proclamation of a democratic Provisional Government based on the authority of the Duma. But the hybrid character of the revolution was once more immediately apparent. Side by side with the Provisional Government, the Petrograd Soviet---the capital had changed its name in 1914---was reconstituted on the model of 1905. 
% 注：清理了 "||"、"appar- ]ent"、"/grad"、"(reconstituted" 等标点乱码和断词。

The revolution of February 1917 brought back to Petrograd, from Siberia and from exile abroad, a host of formerly proscribed revolutionaries. Most of these belonged to one of the two wings---Bolshevik and Menshevik---of the Social-Democratic Workers' Party or to the Social-Revolutionary Party (SRs), and found a ready-made platform in the Petrograd Soviet. The Soviet was in one sense a rival of the Provisional Government set up by the constitutional parties in the old Duma; the phrase ``dual power'' was coined to describe an ambiguous situation. But the initial attitude of the Soviet was less clear-cut. Marx's historical scheme had postulated two distinct and successive revolutions---bourgeois and socialist. Members of the Soviet, with few exceptions, were content to recognize the events of February as the Russian bourgeois revolution which would establish a bourgeois-democratic regime on the western model, and relegated the socialist revolution to a still undetermined date in the future. Cooperation with the Provisional Government was the corollary of this view, which was shared by the first two leading Bolsheviks to return to Petrograd, Kamenev and Stalin. 

Lenin's dramatic arrival in Petrograd at the beginning of April shattered this precarious compromise. Lenin, at first almost alone even among the Bolsheviks, attacked the assumption that the current upheaval in Russia was a bourgeois revolution and nothing more. The situation as it developed after the February revolution confirmed Lenin's view that it could not be confined within bourgeois limits. What followed the collapse of the autocracy was not so much a bifurcation of authority (the ``dual power'') as a total dispersal of authority. The mood of workers and peasants alike, of the vast majority of the population, was one of immense relief at the removal of a monstrous incubus, accompanied by a deep-seated desire to be left to run their own affairs in their own way, and by the conviction that this was somehow practicable and essential. It was a mass movement inspired by a wave of immense enthusiasm and by Utopian visions of the emancipation of mankind from the shackles of a remote and despotic power. It had no use for the western principles of parliamentary democracy and constitutional government proclaimed by the Provisional Government. The notion of centralized authority was tacitly rejected. Local Soviets of workers or peasants sprang up all over Russia. Some towns and districts declared themselves Soviet republics. Factory committees of workers claimed to exercise exclusive authority in their domain. The peasants seized land and divided it among themselves. And everything else was overshadowed by the demand for peace, for an end to the horrors of a bloody and senseless war. 
% 注：清理了大量符号噪点，如 "|lown way"、"[iof"、"“proclaimed"、"lamong"、"fby"、"▼a" 等。

Soldiers' Committees were elected in military units, large and small, from brigades to companies, often demanding the election of officers and challenging their authority. The armies at the front cast off the harsh constraints of military discipline, and slowly began to disintegrate. This all-engulfing movement of revolt against authority seemed to most Bolsheviks a prelude to the fulfilment of their dreams of a new order of society; they had neither the will nor the means to check it. 
% 注：修复了 "...challenging their 1 all-engulfing movement of revolt against authority seemed authority." 这种由于跨行识别导致的严重语序穿插。

When, therefore, Lenin set out to re-define the character of the revolution in his famous ``April theses'', his diagnosis was both perceptive and prescient. He described what had happened as a revolution in transition from a first stage, which had given power to the bourgeoisie, to a second stage, which would transfer power to the workers and the poor peasants. The Provisional Government and the Soviets were not allies but antagonists, representing different classes. The end in view was not a parliamentary republic, but ``a republic of Soviets of Workers', Poor Peasants' and Peasants' Deputies all over the country, growing up from below''. Socialism could not indeed be introduced immediately. But as a first step the Soviets should take control of ``social production and distribution''. Throughout the vicissitudes of the summer of 1917 Lenin gradually secured the adherence of his party followers to this programme. Progress in the Soviets was slower. When an All-Russian Congress of Soviets---the first attempt to create a central Soviet organization with a standing executive committee---met in June, out of more than 800 delegates, the SRs accounted for 285, the Mensheviks for 248, the Bolsheviks only for 105. It was on this occasion that Lenin, in response to a challenge, made the much-derided pronouncement that there was in the Soviet a party ready to take governmental power: the Bolsheviks. As the prestige and authority of the Provisional Government waned, the influence of the Bolsheviks in the factories and in the army grew rapidly; and in July the Provisional Government decided to proceed against them on a charge of conducting subversive propaganda in the army and acting as German agents. Several leaders were arrested. Lenin escaped to Finland, where he carried on a regular correspondence with the party central committee, now working underground in Petrograd. 
% 注：清理了 "[happened"、"Iwhich"、"jwhich" 等异常字符，并将 "social production land distribution" 修正为合乎逻辑的 "social production and distribution"。

It was during this forced withdrawal from the scene of action that Lenin penned one of the most famous, and the most Utopian, of his writings, State and Revolution, a study of Marx's theory of the state. Marx had not only preached the destruction of the bourgeois state by the proletarian revolution, but looked forward, after the victory of the revolution and a transitional period of the dictatorship of the proletariat, to the withering away and eventual extinction of the state. What the proletariat needed at the moment of its victory, Lenin observed, was ``only a state in process of dying away, i.e. so constituted that it will at once begin to die away and cannot help dying away''. The state was always an instrument of class domination and oppression. The classless society of communism and the existence of the state were incompatible. Lenin summed up in an aphorism of his own: ``So long as the state exists, there is no freedom. When freedom exists, there will be no state.'' Lenin was not only a profound student of Marx, but had a sensitive ear for the revolutionary mood of workers and peasants whose enthusiasm was fired by the prospect of escaping from the shackles of a powerful and omnipresent state. State and Revolution was a remarkable synthesis of the teachings of Marx with the aspirations of the untutored masses. The party was scarcely mentioned in its pages. 

By September, after an abortive attempt to seize power by a Right-wing general, Kornilov, the Bolsheviks had won a majority in the Petrograd and Moscow Soviets. Lenin, after some hesitation, revived the slogan ``All power to the Soviets''---a direct challenge to the Provisional Government. In October he returned in disguise to Petrograd to attend a meeting of the party central committee. Under his persuasion the committee decided, with only two dissentients, Zinoviev and Kamenev, to prepare for an immediate seizure of power. The preparation was carried out mainly by a revolutionary military committee which had been set up by the executive committee of the Petrograd Soviet, and which was now firmly in Bolshevik hands. Trotsky, who had joined the Bolsheviks after his return to Petrograd in the summer, played a leading part in planning the operation. On October 25 (Old Style, equivalent to November 7 by the western calendar, which was introduced a few months later) the Red Guard, composed mainly of factory workers, occupied key positions in the city, and advanced on the Winter Palace. It was a bloodless coup. The Provisional Government collapsed without resistance. Some of the ministers were arrested. The Prime Minister, Kerensky, fled abroad. 

The coup had been timed to coincide with the second All-Russian Congress of Soviets of Workers' and Soldiers' Deputies, which opened on the following evening. The Bolsheviks now had a majority---399 out of a total of 649 delegates---and took charge of the proceedings. The congress pronounced the dissolution of the Provisional Government, and the transfer of authority to the Soviets, and unanimously adopted three major decrees, the first two being submitted to it by Lenin. The first was a proclamation, in the name of the ``Workers' and Peasants' Government'', proposing to all the belligerent peoples and governments to enter into negotiations for a ``just and democratic peace'' without annexations or indemnities, and appealing particularly to ``the class-conscious workers of the three most advanced nations of mankind''---Britain, France and Germany---to help to bring the war to an end. 

The second was a decree on land, incorporating a text drawn up by the SRs, which responded to the petty bourgeois aspirations of the peasant rather than to long-term Bolshevik theories of a socialized agriculture. Landlords' ownership of land was abolished without compensation; only the land of ``ordinary peasants and ordinary Cossacks'' was exempt from confiscation. Private ownership of land was abolished for ever. The right of using land was accorded to ``all citizens (without distinction of sex) of the Russian state desiring to cultivate it with their own labour''. Mineral and other subsidiary rights were reserved to the state. The buying, selling and leasing of land, and the employment of hired labour, were prohibited. This was the charter of the small, independent peasant cultivating his plot of land with his own labour and that of his family, and serving primarily their needs. A final settlement of the land question was reserved for the future Constituent Assembly. 
% 注：这一段的错乱最为严重，OCR将两栏文字交叉提取成了诸如 "frights were reserved to the state. The buying, selling and distinction of sex) of the Russian state desiring to cultivate it with their own labour” 等不知所云的句子，现已彻底解开并拼合成原本正常的逻辑和语序。将 ":onfiscation" 修正为 "confiscation"。

The third decree, proposed by Kamenev, who presided at the session, set up a Council of People's Commissars (Sovnarkom) as a Provisional Workers' and Peasants' Government to govern the country under the authority of the All-Russian Congress of Soviets and of its executive committee till the meeting of the Constituent Assembly. 

These pronouncements had several distinctive features. Lenin had ended his speech a few hours earlier in the Petrograd Soviet with the bold words: ``In Russia we must concern ourselves with the building of the proletarian socialist state.'' In the more formal decrees of the Congress of Soviets, the concepts of ``the state'' and of ``socialism'' were kept in the background. In the enthusiasm of victory, when the old state with its attendant evils was being swept away, nobody was eager to face the problem of building a new state. The revolution was international, and took no account of national boundaries. The Workers' and Peasants' Government had no territorial definition or designation; the ultimate extension of its authority could not be foreseen. Socialism was an ideal of the future; Lenin observed, in introducing the decree on peace, that the victory of the workers would ``pave the way to peace and socialism''. But none of the decrees mentioned socialism as the aim or purpose of the revolution; its content, like its extent, was left to the future to decide. 
% 注：去除了 "11 Soviet"、"I 1 ourselves"、"III to peace"、"// socialism" 等明显的噪点数字和斜杠。

Finally, the gesture of deference to the ultimate authority of the Constituent Assembly, which in retrospect seemed oddly illogical, was accepted without objection. Between February and October the Provisional Government and the Soviets had both demanded a constituent assembly---the traditional democratic procedure for the drafting of a new constitution; and the date of November 12 had been fixed for the elections. Lenin did not wish, or did not feel strong enough, to cancel them. As was to be expected in a predominantly rural electorate, the vote gave an absolute majority to the SRs---267 out of 520 deputies; the Bolsheviks had 161, the remainder being made up of a large number of splinter groups. When the deputies met in January 1918, the Workers' and Peasants' Government was firmly established in Petrograd, and was unlikely to abdicate in favour of a body which represented the confused moods of the countryside two months earlier. Bukharin spoke of ``the watershed which at this moment divides this assembly into . . . two irreconcilable camps, camps of principle ... for socialism or against socialism''. The assembly listened to much inconclusive oratory. Late in the night it adjourned; and the government forcibly prevented it from re-assembling. It was a decisive moment. The revolution had turned its back on the conventions of bourgeois democracy. 
% 注：修复了由于版式原因断裂的句子 "(divides this assembly into..."，并删除了 "^ffl"、"^(11" 等乱码字符。

The first consequence of the revolution to impinge on the western world, and to excite horror and indignation, was its withdrawal from the war and its desertion of the Allies at the desperate peak of their struggle with Germany. When this unforgivable betrayal was quickly followed by such measures as the repudiation of the debts of former Russian Governments and the expropriation of owners of land and factories, and when the revolution announced itself as the first stage in a revolution which was destined to sweep over Europe and the world, it was revealed as a fundamental assault on the whole of western capitalist society. But this threat was not taken very seriously. Few people in the west at first imagined that the revolutionary regime in Russia could survive for more than a few days or weeks. The Bolshevik leaders themselves did not believe that they could hold out indefinitely, unless the workers of the capitalist countries came to their aid by rising in revolt against their own governments. 
% 注：清理了 "(flits"、"11 at the"、"1' '"、"y and the"、"[indefinitely" 等识别错误。

This scepticism did not lack plausibility. The writ of the Workers' and Peasants' Government scarcely extended beyond Petrograd and a few other large cities. Even in the Soviets the Bolsheviks did not yet command unanimous support; and it was quite uncertain how far the All-Russian Congress of Soviets---the one sovereign central authority---would be recognized by the local Soviets which had sprung up all over the country, by the factory committees exercising ``workers' control'' in the factories, or by the millions of peasants now flocking back to their homes from the front. Bureaucrats, managers, and technical experts at all levels came out on strike, and refused to serve the new self-styled government. 
% 注：修正 "Gongress" 为 "Congress"。这里原本又发生了一次两栏交织乱序 ("...sprung up all Il“workers’ control” in the factories, or by the millions of over the country... [peasants now flocking...")，已将其完美复位。

The armed forces at the disposal of the regime consisted of a nucleus of a few thousand Red Guards, and some loyal Lettish battalions surviving from the disintegration of the imperial armies which had fought in the war. Within a few weeks of the revolution, Cossack armies pledged to its overthrow were being organized in the regions of the Don, the Kuban, and the Urals. It had been easy for the Bolsheviks to topple the rickety Provisional Government. To substitute themselves for it, to establish effective control over the chaos which had overwhelmed the vast territory of the defunct Russian Empire, and to set up a new order of society geared to the aspirations of the masses of workers and peasants who had seen in the Bolsheviks their saviours and liberators, was a far more formidable and complex task.
% 注：删除了行首乱码 "^" 以及 "=?"。